% BHU PYQ -- Manually transcribed & error-corrected
% Paper: BCF-211 : Corporate Accounting
% Exam: B.Com. (Hons.) F.M.M. Semester III Examination, 2016-17

\documentclass[11pt,a4paper]{article}
\usepackage[a4paper,top=2.5cm,bottom=2.5cm,left=2.8cm,right=2.8cm,headheight=14pt]{geometry}
\usepackage{amsmath,amssymb}
\usepackage[T1]{fontenc}
\usepackage[utf8]{inputenc}
\usepackage{lmodern,microtype,fancyhdr,enumitem,hyperref,lastpage,parskip,booktabs}

\hypersetup{
    colorlinks=true,
    urlcolor=black,
    linkcolor=black,
    pdftitle={BCF-211: Corporate Accounting -- BHU B.Com. (Hons.) FMM Sem III 2016-17}
}

\pagestyle{fancy}
\fancyhf{}
\renewcommand{\headrulewidth}{0.4pt}
\renewcommand{\footrulewidth}{0.4pt}
\fancyhead[L]{\small\textit{Banaras Hindu University}}
\fancyhead[R]{\small\textit{BCF-211: Corporate Accounting}}
\fancyfoot[C]{\small Page \thepage\ of \pageref{LastPage}}

\newlist{parts}{enumerate}{2}
\setlist[parts,1]{label=(\alph*),leftmargin=*,itemsep=5pt,topsep=3pt}
\setlist[parts,2]{label=(\roman*),leftmargin=*,itemsep=3pt,topsep=2pt}
\newcommand{\pts}[1]{\hfill\textit{[#1]}}

\begin{document}

\begin{center}
  {\large\bfseries BANARAS HINDU UNIVERSITY}\\[2pt]
  {\normalsize B.Com. (Hons.) F.M.M. --- Semester III Examination, 2016-17}\\[6pt]
  \rule{0.8\linewidth}{0.5pt}\\[6pt]
  {\large\bfseries Paper No. BCF-211: Corporate Accounting}\\[4pt]
  {\normalsize Time: 3 Hours\hfill Maximum Marks: 70}
\end{center}

\rule{\linewidth}{0.4pt}

\smallskip
\noindent\textit{Instructions: Answer all questions. Marks are noted against each question.}

\bigskip

\noindent\textbf{Question 1.}
Chandra Mills Ltd., Kanpur issues 550, 6\% debentures of Rs.\ 100 each and 4000 equity shares of Rs.\ 10 each to the public. Prakash \& Company has entered into an underwriting agreement with Chandra Mills Ltd.\ for the subscription of the whole of the issue of the public at a commission of 2\% on equity shares and 1\% on debentures. Cash was received for all equity shares and debentures from the public except that of 50 debentures. These were taken up by the underwriters. Underwriters paid the balance due. Pass the necessary journal entries and also prepare the Balance Sheet in the books of Chandra Mills Ltd.\pts{14}

\centerline{\textbf{OR}}

How are shares forfeited? Can forfeited shares be re-issued at a discount? Give Journal entries for the forfeiture and re-issue of shares.\pts{4+4+6}

\medskip\rule{\linewidth}{0.2pt}

\noindent\textbf{Question 2.}
What is amalgamation? Describe the types of amalgamation. Discuss the accounting methods of amalgamation.\pts{2+2+10}

\centerline{\textbf{OR}}

Divya Oil Mills Ltd.\ go into voluntary liquidation. A new company Ganesh Oil Mills Ltd.\ was formed to acquire the assets of Divya Oil Mills Ltd.\ on 31st March, 2016. Assets and liabilities of Divya Oil Mills Ltd.\ were as follows: Works and other properties Rs.\ 3,60,000; current assets Rs.\ 40,000; current liabilities Rs.\ 80,000. Its paid up capital was Rs.\ 4,00,000. The assets were sold to the new company for Rs.\ 2,88,000, Rs.\ 2,00,000 payable in shares of Rs.\ 1 each, credited with 75 paisa per share paid up and Rs.\ 88,000 in cash. Divya Oil Mills Ltd.\ paid the liabilities and the cost of liquidation. Pass the journal entries in the books of Divya Oil Mills Ltd.\pts{14}

\medskip\rule{\linewidth}{0.2pt}

\noindent\textbf{Question 3.}
What is meant by Internal Reconstruction? What accounting entries are passed in the books of companies at the time of Internal Reconstruction?\pts{4+10}

\centerline{\textbf{OR}}

On the reconstruction of a Company the following conditions were agreed upon:
The Shareholders to receive in lieu of their present holding (i.e.\ 25,000 shares of Rs.\ 10 each) the following:
\begin{parts}
  \item Fully paid equity shares to 2/5 of their holding.
  \item 5\% preference shares, fully paid to the extent of 1/5 of the above new equity shares.
  \item Rs.\ 30,000, 6\% second debentures.
\end{parts}
An issue of Rs.\ 25,000, 5\% first debentures was made and allotted, payment for the same having been received in cash. The Goodwill, which stood at Rs.\ 1,50,000 was written down to Rs.\ 75,000. The land and building, which stood at Rs.\ 50,000 was written down to Rs.\ 37,500. The plant \& machinery, which stood at Rs.\ 75,000 were written-down to Rs.\ 62,500. Give the journal entries in the books of the company on the basis of the above description.\pts{14}

\medskip\rule{\linewidth}{0.2pt}

\noindent\textbf{Question 4.}
Describe preferential creditors according to the Companies Act. Give a specimen of the Liquidator's final statement of accounts with imaginary figures.\pts{5+9}

\centerline{\textbf{OR}}

A Company Ltd.\ went into liquidation with the following liabilities:
\begin{parts}
  \item Secured creditors Rs.\ 20,000 (securities realized Rs.\ 25,000)
  \item Preferential creditors Rs.\ 6,000
  \item Unsecured creditors Rs.\ 30,500
\end{parts}
Liquidation Expenses Rs.\ 252.

Liquidator's remuneration: 3\% on the amount realized and $1\frac{1}{2}\%$ on the amount distributed to unsecured creditors. Rs.\ 26,000 were realized from various assets and this amount does not include those securities which are with the secured creditors. Prepare the Liquidator's Final statement of account.

The liquidator is entitled to remuneration on the amount realized on all assets including those which are with the secured creditors.\pts{14}

\medskip\rule{\linewidth}{0.2pt}

\noindent\textbf{Question 5.}
Explain the following:\pts{$3\frac{1}{2} \times 4$}
\begin{parts}
  \item Minority shareholder's interest.
  \item Amount of Goodwill for Consolidated Balance Sheet.
  \item Consolidated Balance Sheet.
  \item Cost of control.
\end{parts}

\centerline{\textbf{OR}}

What is a consolidated Balance Sheet? How is it prepared? Give its specimen.\pts{4+5+5}

\vfill
\rule{\linewidth}{0.4pt}
\begin{center}\small
\href{https://bhu-exam.vercel.app/}{Click here to visit our website}\\[4pt]
\href{https://me-aryan.vercel.app/}{Click here to visit our contributor}\\[4pt]
\href{https://quashan-me.vercel.app/}{Click here to visit our contributor}
\end{center}

\end{document}
